\section{Inequalities: Markov, Chebyshev, Jensen, and norm bounds}\label{sec:inequalities}

\subsection{Markov and Chebyshev}

\begin{theorem}[Markov inequality]\label{thm:markov}
Let \(X\ge0\) pointwise and let \(a>0\) be rational.  Then
\[
 \Pp(X\ge a)\le\frac{\E[X]}{a}.
\]
Equivalently, \(a\Pp(X\ge a)\le\E[X]\).
\end{theorem}

\begin{proof}
At each outcome,
\[
 a\ind_{X\ge a}\le X.
\]
Indeed, outside the event the left side is zero, and on the event it is \(a\),
which is at most \(X\).  Take expectations, use monotonicity and the indicator
formula, and divide by the positive \(a\).  The object-language target is the
cross-multiplied inequality, so no division function is required.
\end{proof}

\begin{corollary}[Power Markov]\label{cor:power-markov}
For natural \(r>0\), rational \(a>0\), and any \(X\),
\[
 \Pp(\abs X\ge a)\le\frac{\E[\abs X^r]}{a^r}.
\]
\end{corollary}

\begin{proof}
Apply \Cref{thm:markov} to the nonnegative variable \(\abs X^r\).  The events
\(\abs X^r\ge a^r\) and \(\abs X\ge a\) coincide because exponentiation by a
positive natural is order preserving on nonnegative rationals.  This last fact
is proved by induction on \(r\).
\end{proof}

\begin{theorem}[Chebyshev inequality]\label{thm:chebyshev}
For \(a>0\),
\[
 \Pp(\abs{X-\E[X]}\ge a)\le\frac{\Var(X)}{a^2}.
\]
\end{theorem}

\begin{proof}
Apply power Markov with \(r=2\) to \(X-\E[X]\).  Its second absolute moment is
exactly the variance.
\end{proof}

\begin{corollary}[At least three quarters within two standard deviations]
\label{cor:three-quarters}
If a positive rational \(s\) satisfies \(s^2=\Var(X)\), then
\[
 \Pp(\abs{X-\E[X]}<2s)\ge\frac34.
\]
The same denominator-free statement is
\(4\Var(X)\Pp(\abs{X-\E[X]}\ge2s)\le\Var(X)\), with the zero-variance case
handled separately.
\end{corollary}

\begin{proof}
Chebyshev gives probability at most \(1/4\) to the complementary event.
Subtract from one using the finite event identity.  A square root need not be
constructed internally if the theorem is stated with an input \(s\) and the
certificate \(s^2=\Var(X)\).
\end{proof}

\subsection{Constructive Jensen by a supporting-line certificate}

A bare classical definition of convexity quantifies over real coefficients.
For \(\ipra\), the most useful finite formulation is a checkable supporting
line at the particular mean under consideration.

\begin{definition}[Jensen certificate]\label{def:jensen-cert}
Let \(\phi\) be a primitive-recursive rational evaluator on the finite list of
values \(X_i\) and on \(\mu=\E[X]\).  A Jensen certificate at \(\mu\) is a
rational \(s\) such that for every \(i<m\),
\[
 \phi(X_i)\ge \phi(\mu)+s(X_i-\mu).
\]
\end{definition}

\begin{theorem}[Certified finite Jensen inequality]\label{thm:jensen}
If \(s\) is a Jensen certificate for \(\phi\) at \(\mu=\E[X]\), then
\[
 \phi(\E[X])\le\E[\phi(X)].
\]
\end{theorem}

\begin{proof}
Multiply the certificate inequality by \(w_i\), sum, and divide by \(W\):
\[
 \E[\phi(X)]
 \ge \phi(\mu)+s\E[X-\mu]
 =\phi(\mu).
\]
The vanishing of the final expectation is linearity and the constant rule.
All certificate checks are finite rational comparisons.
\end{proof}

\begin{corollary}[Square and absolute-value Jensen]\label{cor:jensen-basic}
For every finite rational \(X\),
\[
 \E[X]^2\le\E[X^2],\qquad
 \abs{\E[X]}\le\E[\abs X].
\]
\end{corollary}

\begin{proof}
For \(\phi(x)=x^2\), choose \(s=2\mu\); the certificate difference is
\[
 x^2-\mu^2-2\mu(x-\mu)=(x-\mu)^2\ge0.
\]
For \(\phi(x)=\abs x\), decide the sign of \(\mu\).  If \(\mu\ge0\), choose
\(s=1\) and use \(\abs x\ge x\); if \(\mu<0\), choose \(s=-1\) and use
\(\abs x\ge-x\).
\end{proof}

\begin{proposition}[Piecewise-linear rational Jensen]\label{prop:polygonal-jensen}
Suppose \(\phi\) is given by a finite rational convex polygon: rational
breakpoints \(b_0<\cdots<b_r\) and nondecreasing rational slopes
\(s_0\le\cdots\le s_{r-1}\).  Then a Jensen certificate at every rational
\(\mu\) in the domain is primitive-recursively computable, and
\Cref{thm:jensen} applies.
\end{proposition}

\begin{proof}
Bounded search locates the interval containing \(\mu\).  Choose its slope, or
at a breakpoint either adjacent slope.  On each side of \(\mu\), telescope the
polygonal increments.  Nondecreasing slopes show that the graph lies above the
chosen affine line.  The argument is a finite sum of products of nonnegative
rational segment lengths and slope differences.
\end{proof}

\begin{remark}[Smooth convex functions]
For a smooth function represented by regular rational approximants, the same
proof works with an approximate supporting slope and a certified error term.
At precision \(p\), use a rational slope whose supporting inequality is within
\(2^{-p-2}\), apply finite Jensen, and absorb the endpoint errors.  Thus the
usual Jensen theorem for effectively convex functions is an effective analytic
schema, not a new logical principle.
\end{remark}

\subsection{Cauchy--Schwarz as a finite sum of squares}

\begin{lemma}[Weighted Lagrange identity]\label{lem:lagrange}
For rational lists \(u_i,v_i\) and natural weights \(w_i\), let
\[
 A=\sum_iw_iu_i^2,\quad B=\sum_iw_iv_i^2,\quad C=\sum_iw_iu_iv_i.
\]
Then
\[
 AB-C^2
 =\sum_{0\le i<j<m}w_iw_j(u_iv_j-u_jv_i)^2.
\]
\end{lemma}

\begin{proof}
Expand the left side as a double sum:
\[
 AB-C^2=\sum_{i,j}w_iw_j(u_i^2v_j^2-u_iv_iu_jv_j).
\]
The diagonal terms vanish.  Pair the terms \((i,j)\) and \((j,i)\) for
\(i<j\); their sum is
\[
 w_iw_j(u_i^2v_j^2+u_j^2v_i^2-2u_iv_iu_jv_j)
 =w_iw_j(u_iv_j-u_jv_i)^2.
\]
Finite pairing and Fubini are justified by \Cref{cor:finite-fubini}.  The
identity is polynomial and can also be discharged by the ProveIt ring tool
after the double-sum lemmas have been instantiated.
\end{proof}

\begin{theorem}[Cauchy--Schwarz]\label{thm:cauchy-schwarz}
For finite rational observables,
\[
 \E[XY]^2\le\E[X^2]\E[Y^2].
\]
Consequently
\[
 \abs{\Cov(X,Y)}^2\le\Var(X)\Var(Y).
\]
\end{theorem}

\begin{proof}
Apply \Cref{lem:lagrange} with \(u_i=X_i\) and \(v_i=Y_i\).  The right side is
nonnegative, so \(C^2\le AB\).  Divide by \(W^2>0\).  For covariance, apply
the first inequality to \(X-\E[X]\) and \(Y-\E[Y]\).
\end{proof}

\begin{corollary}[Equality criterion]\label{cor:cs-equality}
Equality holds in Cauchy--Schwarz precisely when
\[
 w_iw_j(X_iY_j-X_jY_i)^2=0
\]
for every \(i<j\).  In particular, on the positive-weight support the two
vectors are rationally proportional unless both vanish in the relevant
coordinates.
\end{corollary}

\begin{proof}
The difference in \Cref{lem:lagrange} is a finite sum of nonnegative rational
terms.  It is zero exactly when every term is zero, by zero-sum elimination.
The proportionality reading follows by selecting any positive-weight index
where one coordinate is nonzero and cross-multiplying.
\end{proof}

\subsection{Holder, Lyapunov, and Minkowski}

The article lists several norm inequalities.  For integer exponents one can
prove them in exact algebraic form; for arbitrary computable real exponents we
state the regular-real version with root and power moduli.

\begin{theorem}[Finite Holder with effective rational exponents]
\label{thm:holder}
Let \(p>1\) and \(q>1\) be fixed positive rational exponents satisfying
\(1/p+1/q=1\).  Assume the presentation includes regular evaluators for the
rational powers and roots, together with a checked Young-inequality certificate
at the requested precision.  Then
\[
 \E[\abs{XY}]
 \le \E[\abs X^p]^{1/p}\E[\abs Y^q]^{1/q}.
\]
The special case \(p=q=2\) is the exact rational theorem
\Cref{thm:cauchy-schwarz}.  For general rational conjugates, the statement is
an effective regular-real schema.  Its finite object-language core says that
rational upper bounds \(A,B>0\) satisfying
\(A^p\ge\E\abs X^p\) and \(B^q\ge\E\abs Y^q\), to certified precision, imply
\(\E\abs{XY}\le AB\) up to the allocated dyadic error.
\end{theorem}

\begin{proof}
Normalize by positive rational upper bounds \(A,B\).  Young's inequality
\(uv\le u^p/p+v^q/q\) for nonnegative inputs is verified at finite precision by
the supplied power certificates; for fixed integer powers it can instead be
reduced to a polynomial factorization after clearing denominators.  Apply it to
\(u=\abs X/A\) and \(v=\abs Y/B\), multiply by \(AB\), and take expectations.
The normalized moments are at most one, yielding the claimed bound, with the
certificate errors absorbed by guard bits.  Regular roots are found by
bisection, so their interval widths provide a primitive-recursive output
modulus.
\end{proof}

\begin{corollary}[Lyapunov moment monotonicity]\label{cor:lyapunov}
For \(0<r<s\) given by fixed positive rational exponents and a finite rational
observable,
\[
 \E[\abs X^r]^{1/r}\le\E[\abs X^s]^{1/s}.
\]
\end{corollary}

\begin{proof}
Apply Holder to \(\abs X^r\cdot1\) with conjugate exponents
\(s/r\) and \(s/(s-r)\), then use \(\E[1]=1\).  In the denominator-cleared
form, all roots appear only through certified rational enclosures.
\end{proof}

\begin{theorem}[Finite Minkowski]\label{thm:minkowski}
For fixed computable \(p\ge1\),
\[
 \E[\abs{X+Y}^p]^{1/p}
 \le \E[\abs X^p]^{1/p}+\E[\abs Y^p]^{1/p}.
\]
For \(p=1\) the theorem is exact rational algebra; for \(p=2\) it follows
from Cauchy--Schwarz; for general fixed rational \(p\) it is an effective
regular-real schema based on Holder.
\end{theorem}

\begin{proof}
For \(p=1\), use the pointwise triangle inequality and monotonicity.  For
\(p>1\), put \(q=p/(p-1)\).  The pointwise inequality
\[
 \abs{X+Y}^p
 \le \abs X\abs{X+Y}^{p-1}+\abs Y\abs{X+Y}^{p-1}
\]
followed by expectation and Holder gives
\[
 M^p\le (A+B)M^{p-1},
\]
where \(M=\norm{X+Y}_p\), \(A=\norm X_p\), and \(B=\norm Y_p\).
If \(M=0\) the claim is immediate; otherwise cancel the positive
\(M^{p-1}\).  Constructively, cancellation is performed on rational interval
bounds separated from zero; when separation is unavailable, the target error
already dominates \(M\), so the approximate inequality closes.  This standard
case split is decidable at each finite precision.
\end{proof}
